\section{Endüstri Deneyimi}

\begin{twocolentry}{
    \textit{Şub 2022 – Eki 2022}\\
    \small Uzaktan}
    \normalsize \textbf{Yazılım Geliştirme Stajyeri}, \small {Mainflux}
    \small  
    \begin{itemize}
        \item InfluxDB'yi sürüm 2'ye yükseltti ve tüm sorguları InfluxQL'den FluxQL'e taşıdı. Sorgu karşılaştırıcılarındaki sorunları giderdi ve yeni karşılaştırma operatörleri ekledi.\\
        \footnotesize \textbf{Araçlar:} Golang, Docker, nGinx, Go-Kit, InfluxDB, Mikroservis Mimarisi
    \end{itemize}
\end{twocolentry}

\vspace{\betweenentryspace}

\begin{twocolentry}{
    \textit{Ağu 2021 – Eki 2021}\\
    \small İstanbul, Türkiye}
    \normalsize \textbf{Yazılım Geliştirme Stajyeri}, \small {Tübitak Bilgem YTE (Türkiye Ulusal Araştırma Enstitüsü)}
    \small  
    \begin{itemize}
        \item Bir etkinlik yönetimi ve pazarlama platformu için ön uç ve arka uç geliştirdi.\\
        \footnotesize \textbf{Araçlar:} Spring Boot, React.js, PostgreSQL, Spring Security
    \end{itemize}
    
\end{twocolentry}

\vspace{\betweenentryspace}

\begin{twocolentry}{
    \textit{Haz 2020 – Ağu 2020}\\
    \small İstanbul, Türkiye}
    \normalsize \textbf{Veri Bilimi Stajyeri}, \small {Faraday Networks}
    \small  
    \begin{itemize}
        \item Müşteri Wi-Fi oturumlarının zaman serisi verilerinde anormallik tespiti yapmak için makine öğrenmesi modelleri kullanan bir araç geliştirdi.\\
        \footnotesize \textbf{Araçlar:} PostgreSQL, Python, Facebook Prophet
    \end{itemize}
    
\end{twocolentry}
